\bigheading{Supplementary Discussion 1: Temporal hydrogen regulation}
\label{subsec:gh_constraint_effects}

\textbf{Electricity supply and demand without temporal hydrogen regulation}. 
Supplementary Fig. \ref{fig:balances-ac-co2l20-120-noreg} shows the effect of increasing domestic \co mitigation at constant (1 TWh) hydrogen exports without temporal matching. Since there are no substantial hydrogen exports, the difference of temporal matching (displayed in Figure \ref{fig:balances-ac-co2l20-120}) is only marginal. Supplementary Fig. \ref{fig:balances-ac-0exp-120-noreg} displays the hydrogen export ramp up from 1--120 TWh/a at 0\% domestic \co mitigation without temporal regulation. If the no temporal regulation is applied, the dispatch of coal and gas generators increase, which is in contrast to the hydrogen export ramp up with hourly matching, displayed in Figure \ref{fig:balances-ac-0exp-120}.

\begin{figure*}[h!]
    \centering
    \begin{subfigure}[b]{0.49\linewidth}
        \centering
        \includegraphics[trim={0cm 0cm 0cm 1cm}, clip, width=\linewidth]{../../workflow/subworkflows/pypsa-earth-sec/results/\runsensnogreenhy/0exp-only/120/graphs/balances-AC.pdf}
        \caption{Fixed (1 TWh) export and 0--100\% domestic \co mitigation}
        \label{fig:balances-ac-0exp-120-noreg}
    \end{subfigure}
    \hfill
    \begin{subfigure}[b]{0.49\linewidth}
        \centering
        \includegraphics[trim={0cm 0cm 0cm 1cm}, clip, width=\linewidth]{../../workflow/subworkflows/pypsa-earth-sec/results/\runsensnogreenhy/co2l20-only/120/graphs/balances-AC.pdf}
        \caption{Fixed (0\%) domestic \co mitigation and 1--120 TWh/a export}
        \label{fig:balances-ac-co2l20-120-noreg}
    \end{subfigure}
    \hfill
    \caption{Supplement to Figure \ref{fig:balances-ac} without temporal hydrogen regulation. Electricity supply and demand at fixed export levels and increasing domestic \co mitigation (see Supplementary Fig. \ref{fig:balances-ac-0exp-120-noreg}) and vice versa (see Supplementary Fig. \ref{fig:balances-ac-co2l20-120-noreg}). Increasing domestic \co mitigation first phases out carbon-intensive coal generation in favor of CCGT, at medium to high domestic \co mitigation the electricity system is fully renewable supported by flexibility through Vehicle-to-Grid (V2G) and sector coupling. Increasing electricity demands include Battery Electric Vehicles (BEV) and hydrogen generation for other sectors. At increasing hydrogen exports the additional electricity required for hydrogen electrolysis is covered by onshore wind and solar PV, as well as coal and gas power plants since no temporal hydrogen regulation is in place.}
    \label{fig:balances-ac-noreg}
\end{figure*}


\textbf{Effects of temporal matching in a high export and low \co mitigation scenario on total system costs}. Supplementary Fig. \ref{fig:tsc-120-0} displays the total system costs at 120 TWh/a export and 0\% domestic \co mitigation. Stricter temporal hydrogen regulation mainly increases the total CAPEX of additional solar PV, electrolysis and hydrogen storage. In return, the OPEX of fossil generation (mainly gas), decreases. The large share of oil OPEX is independent of temporal hydrogen regulation, since these costs are mainly linked to combustion engine cars with demands independent of temporal hydrogen regulation.

\begin{figure*}[h]
    \centering
    \includegraphics[trim={0cm 0cm 0cm 0.65cm}, clip, width=0.6\linewidth]{../../results/gh/abs_sin_tech_120_Co2L2.0.pdf}
    \caption{Total system costs at 120 TWh/a export and 0\% domestic \co mitigation. Stricter temporal hydrogen regulation mainly increases the total CAPEX of additional solar PV, electrolysis and hydrogen storage. In return, the OPEX of fossil generation (mainly gas), decreases. The large share of oil OPEX is independent of temporal hydrogen regulation, since these costs are mainly linked to combustion engine cars with demands independent of temporal hydrogen regulation.}
    \label{fig:tsc-120-0}
\end{figure*}


\textbf{Cost for domestic electricity consumers and hydrogen exporters without hydrogen regulation}.
Supplementary Fig. \ref{fig:expenses_default_120-noreg} shows, that increasing hydrogen exports quickly in a system that is still dominated by fossil fuels substantially raises costs for domestic electricity consumers, if green hydrogen production is not regulated. At low hydrogen exports (1-20 TWh), an increase of domestic \co mitigation increases the cost for hydrogen exporters.


\begin{figure*}[h!]
    \centering
    \begin{subfigure}[b]{0.49\linewidth}
        \centering
        \includegraphics[width=\linewidth]{../results/\runsensnogreenhy/graphics/integrated_comp/contour_exp_AC_exclu_H2 El_all_20_filterTrue_nTrue_exp120.pdf}
        \caption{Impact of hydrogen exports on domestic interests:        
        Relative cost of electricity for domestic consumers (normalized to 1 TWh/a hydrogen export)}
        \label{fig:expense_ac_120-noreg}
    \end{subfigure}
    \hfill
    \begin{subfigure}[b]{0.49\linewidth}
        \centering
        \includegraphics[width=\linewidth]{../results/\runsensnogreenhy/graphics/integrated_comp/contour_exp_H2_False_False_exportonly_20_filterTrue_nTrue_exp120.pdf}
        \caption{
        Impact of domestic \co mitigation on hydrogen exporter interests:   
        Relative cost of hydrogen for exporters (normalized to 0\% domestic \co mitigation)}
        \label{fig:expense_h2_120-noreg}
    \end{subfigure}
    \hfill
    \caption{  
    Supplement to Figure \ref{fig:expenses_default_120} without temporal hydrogen regulation. Supplementary Fig. \ref{fig:expense_ac_120-noreg} shows the effect of hydrogen exports on domestic electricity cost. Therefore, the effect of hydrogen exports is isolated by normalizing the costs to 1 TWh/a hydrogen export in each column. This approach allows for a vertical interpretation in Supplementary Fig. \ref{fig:expense_ac_120-noreg} only, displaying the effects of hydrogen exports (vertical) at a certain domestic \co mitigation level.
    Supplementary Fig. \ref{fig:expense_h2_120-noreg} shows the effect of domestic \co mitigation on hydrogen export cost. Therefore, the effect of domestic \co mitigation is isolated by normalizing the costs to 0\% \co mitigation in each row. This approach allows for a horizontal interpretation in Supplementary Fig. \ref{fig:expense_h2_120-noreg} only, displaying the effects of domestic \co mitigation (horizontal) at a certain hydrogen export quantity.
    Increasing hydrogen exports quickly in a system that is still dominated by fossil fuels substantially raises costs for domestic electricity consumers, if green hydrogen production is not regulated.
    At low hydrogen exports (1-20 TWh), an increase of domestic \co mitigation increases the cost for hydrogen exporters.
    }
    \label{fig:expenses_default_120-noreg}
\end{figure*}



\clearpage